\chapter[Introduction]{Introduction}
\label{cp:introduction}

{
    \parindent0pt

    Sub-Saharan Africa faces a significant waste management crisis. The consequences of improperly managed waste are severe, leading to environmental pollution, the spread of disease, and the degradation of community spaces. Formal collection systems are often overburdened, leaving communities to bear the brunt of the problem. However, this challenge also creates an opportunity for innovation.

    \vspace{.935em}
    Ayang, named from the Ibibio word for a sweeping broom, is a web-based platform that transforms waste collection into a rewarding, community-building activity. It uses gamification to motivate user participation , a method proven to drive pro-environmental behavior.

    \vspace{.935em}
    The technical core of Ayang is Google's Gemini Vision model. Users photograph a littered site before and after cleaning. The AI assesses the trash and verifies the cleanup, ensuring the integrity of the system. Points are awarded based on the estimated weight of the collected trash, directly incentivizing impactful cleanups. This project aligns with the DLI 2025 Community Challenge by offering a practical, AI-driven solution to a real-world African problem, designed to improve local environments and strengthen community bonds.

    \vspace{.935em}
    This document provides a condensed overview of the Ayang project. For a comprehensive review of our methodology, data, and technical architecture, please access the full report.

    \vspace{.935em}
    Read the full detailed report: \href{https://drive.google.com/file/d/1dQtCNl9HtdUR7oaR-Z_yMTsNeOd9Zuac/view?usp=sharing}{link}

    Try the Ayang web app: \url{https://ayang-ai.vercel.app/}
}